We now present the proof sketch of our main result on the  \cliff{} function, together with 
all supporting lemmata.  All proofs are deferred to the appendix; the 
statements here are mathematically identical to those in the full technical 
development, but streamlined for readability.  The argument proceeds in two 
major stages: (i) obtaining a canonical entry point onto the second slope 
immediately after an ageing event at the cliff-top; and (ii) showing that from 
this configuration the algorithm reaches the optimum without dangerous increases in mutation potential.


\begin{theorem}[Main Runtime Bound] \label{th:expoHD-cliff}
	Let $d < n(\tfrac14-\varepsilon)$ for fixed $\varepsilon>0$, and assume IPH 
	with linear mutation potential  $M(x)=\max\{1,\lfloor c\,\mathrm{HD}(x,\textcolor{purple}{x^\top})\rfloor\}$ 
	for some 
	\[
	c < \frac{1+4\varepsilon}{3-4\varepsilon}.
	\]
	If the ageing threshold satisfies $\tau=\Omega(n^{1+\varepsilon'})$ for some 
	$\varepsilon'>0$, then the $(1+1)$~Opt-IA$_{\textsc{FirstBest}}$ optimises \cliff{} in expected
	\[
	O\!\left(\frac{n^{5}\log(n/d)}{d^{3}} + n^{2}\log d\right)
	\]
	fitness evaluations.
\end{theorem}

To establish this result, we first analyse how the algorithm reaches a highly 
structured ``canonical'' state shortly after ageing at the cliff-top.
At the cliff-top the algorithm cannot improve except by sampling the global 
optimum, and so ageing will eventually be triggered.  The key insight is that 
the iteration of ageing followed by two additional \textcolor{blue}{fitness improvements} can place the 
algorithm into a configuration with \textcolor{purple}{$n-d+2$ 1-bits and the minimal Hamming distance of two to the $best$ solution which had $n+d$ 1-bits.}  The following lemma formalises this.

\begin{lemma}[Canonical Start After Cliff-Top Ageing]\label{lem:can-start-omega-d3}
	Suppose ageing is triggered at the cliff-top (zero-count~$d$).  
	Under IPH with $c\in(0,1)$, the reference-update rule, and ageing resets only 
	on strict improvements, the algorithm reaches within at most three steps the 
	configuration
	\[
	r_0=2, \qquad \mathrm{HD}(r_0)=2,
	\]
	with probability at least 
	\[
	\frac{3}{256}\left(\frac{d}{n}\right)^3.
	\]
\end{lemma}

The canonical condition $(r_0,\mathrm{HD}(r_0)) =
(2,2)$ means that we begin the drift analysis two successful improvements into
the second slope, at a point differing from the reference in exactly two
bit positions.

Repeated ascents to the cliff-top therefore form a geometric waiting-time 
process.  Since each ascent takes $\Theta(n^2\log(n/d))$ evaluations, the 
total time to obtain the canonical start is \new{provided by the following corollary}.

\begin{corollary}[Expected Time to Canonical Start]
	The expected evaluation count until the first occurrence of 
	$(r_0,\mathrm{HD}(r_0)) = (2,2)$ is 
	\[
	O\!\left(\frac{n^{5}\log(n/d)}{d^{3}}\right).
	\]
\end{corollary}


From the canonical starting point, the key difficulty is preventing the mutation
potential from increasing too much.  With only few zero-bits remaining, successful
steps on the second slope require flipping one of these rare zeros, whereas each
$1\!\to0$ flip increases the Hamming distance and thus the mutation potential.
If this potential grows too large, hypermutations become long enough that, under
FCM$(\ge)$, they are increasingly likely to wander back to the cliff-top; once
there, the stopping rule causes the process to lock onto the plateau and erase
all progress.  Hence the analysis must show that, with constant probability, the
mutation potential stays small throughout the entire second-slope climb.  This
requires tightly controlling how many detrimental flips a hypermutation can
accumulate before recovering to parent-level fitness.  The next lemma provides
the crucial starting point: an exponential tail bound for this damaging prefix.

\begin{lemma}[Exponential Tail for Catch-Up Length] \label{lem:xi-tail}
	Let the current solution on the second slope have $z\le d$ zero-bits and 
	define $p=z/n$.  
	Let $\xi_z$ be the number of $1\to0$ flips before fitness first returns to 
	$\ge$ the parent's.  
	Then for all $k\ge1$,
	\[
	\Pr(\xi_z\ge k)\le (2p)^k.
	\]
	Consequently, the probability that Hamming distance increases by at least 
	$2k$ in a single hypermutation is also at most $(2p)^k$.
\end{lemma}

A whole \emph{phase} consists of all attempts at the same zero-count until the 
first success.  
Using the tail bound above, we obtain uniform exponential moment bounds for 
entire phases.

\begin{lemma}[Phase Moment-Generating Bound] \label{lem:phase-mgf}
	Consider a phase at zero-count $z$ and let $Y_z$ denote the total 
	Hamming-distance increase from all failed attempts in the phase.  
	Let $\rho = 2d/n$.  
	For all $0<\lambda<\tfrac12\ln\frac{1}{2\rho}$,
	\[
	\mathbb{E}[e^{\lambda Y_z}] 
	\;\le\; 
	K(\lambda)
	:= 
	\frac{1-\rho}{1-\rho e^{2\lambda}}.
	\]
\end{lemma}

To track the safety of mutation potential across phases, we introduce a 
linear \emph{buffer} \new{$B_r = \theta r - s_0 - E_r$} capturing the invariant $c\,\mathrm{HD}(r)\le r$.

\begin{lemma}[Buffer Process and Safety Criterion] \label{lem:buffer}
	Let successful phases be indexed by $r=r_0,r_0+1,\dots$ and define
	\[
	s_0=\mathrm{HD}(r_0)-r_0, \qquad 
	\theta=\frac{1}{c}-1, \qquad
	E_r=\sum_{j=r_0+1}^{r} Y_{z_j},
	\]
	and $B_r = \theta r - s_0 - E_r$.
	Then for all $r\ge r_0$:
	\begin{enumerate}
		\item 
		$B_r\ge0$ if and only if $c\cdot\mathrm{HD}(r)\le r$;
		\item 
		$B_{r+1}-B_r=\theta-Y_{z_{r+1}}$;
		\item 
		$B_{r_0}=\theta r_0 - s_0$.
	\end{enumerate}
\end{lemma}

Combining the phase MGF with the buffer update yields an exponential drift 
envelope suitable for applying \textcolor{purple}{Ville's inequality for non-negative supermartingales (Theorem~\ref{})}.

\begin{lemma}[Exponential Drift Envelope] \label{lem:drift-envelope-global}
	Let $\rho=2d/n$.  
	For any $0<\lambda<\tfrac12\ln\tfrac{1}{2\rho}$ define
	\[
	K(\lambda)=\frac{1-\rho}{1-\rho e^{2\lambda}}, 
	\qquad 
	\alpha=e^{-\lambda\theta}K(\lambda).
	\]
	Then for all $t\ge0$,
	\[
	\mathbb{E}[e^{-\lambda(B_{t+1}-B_t)}\mid\mathcal{F}_t]\le \alpha.
	\]
\end{lemma}

\new{We now apply Ville's inequality to turn such a drift envelope into a bound on barrier crossing.}


\begin{corollary}[Buffer Barrier Bound] \label{cor:ville-exponential}
	
	If $B_0=b>0$ and there exist $\lambda>0$ and $\alpha<1$ with
	\[
	\mathbb{E}[e^{-\lambda(B_{t+1}-B_t)}\mid\mathcal{F}_t]\le\alpha,
	\]
	then
	\[
	\Pr\!\left(\inf_{t\ge0} B_t \le 0\right) \le e^{-\lambda b}.
	\]
\end{corollary}

To guarantee the drift negativity required above, we identify the admissible 
range of the mutation-potential factor~$c$.

\begin{lemma}[Per-Instance $c$-Threshold] \label{lem:c-threshold-instance}
	
	Let $\rho=2d/n$ and $\theta=\frac{1}{c}-1$.  
	Define
	\[
	c_\star(n,d)=\frac{1-\rho}{1+\rho}.
	\]
	If $c<c_\star(n,d)$, then there exists $\lambda>0$ such that
	\[
	e^{-\lambda\theta}K(\lambda) < 1,
	\]
	and hence the buffer satisfies the Ville barrier bound.
\end{lemma}

A uniform instantiation over all permitted~$d$ shows that the feasible~$c$ 
range has a nontrivial lower bound.

\begin{corollary}[Uniform Lower Bound for $c$]
	If $d\le (\tfrac14-\varepsilon)n$, then
	\[
	c_\star(n,d) \;\ge\; \frac{1+4\varepsilon}{3-4\varepsilon} \;\ge\; \tfrac13.
	\]
\end{corollary}
\label{cor:c-floor}


Finally, combining the threshold on $c$ with the initial buffer value at 
$(r_0,\mathrm{HD}(r_0))=(2,2)$ yields a uniform constant bound on the 
probability of avoiding a jump-back.

\begin{lemma}[Constant Safety Probability] \label{lem:param-constant-c-general}
	If $c<c_\star(n,d)$ and 
	$b=B_{r_0}=\frac{r_0}{c}-\mathrm{HD}(r_0)>0$, 
	then 
	\[
	\Pr\!\left(\inf B_t \le 0\right) \le e^{-\lambda b}
	\]
	for some $\lambda>0$ depending only on $c$ and $d/n$.  
	For fixed~$\varepsilon$, this yields a uniform constant $<1$ independent of 
	$n,d$.
\end{lemma}

Once the mutation potential remains safe throughout, progress on the second 
slope is monotone and its runtime is easily bounded.



\begin{lemma}[Second-Slope Climb Time] \label{prop:climb-time}
	
	If $B_t>0$ for all $t$, then the second-slope ascent completes in expected
	\[
	\Theta(n\log d)
	\]
	iterations.  
	If $\tau=\Omega(n^{1+\varepsilon})$, then $\Theta(n\log d)=o(\tau)$.
\end{lemma}


The expected time to reach the canonical start is 
$O(n^{5}\log(n/d)/d^{3})$.  
Once there, with constant probability the buffer remains positive for all 
phases; conditioned on this event, the second-slope ascent costs 
$O(n^{2}\log d)$ evaluations.  
Since the success probability is constant, this remains the unconditional 
expectation.  
Summing the two contributions establishes the main runtime bound.
